\begin{trapbox}
\begin{itemize}
    \item \textbf{符号混淆}：德摩根公式中的 $\cap$ 和 $\cup$ 易混淆。记忆技巧：将补集符号看作一个"长条"，当它覆盖整个括号时，括号内的运算要变号（交变并，并变交），并且补集符号分配到每个集合上。
    \item \textbf{全集变化}：德摩根公式依赖于全集 $U$。如果全集不同，补集定义不同，公式仍成立但具体元素会变。做题时要注意全集是否一致。
    \item \textbf{逻辑等价}：德摩根公式在逻辑中对应 $\neg (P \land Q) \equiv (\neg P) \lor (\neg Q)$ 和 $\neg (P \lor Q) \equiv (\neg P) \land (\neg Q)$。在集合论中应用时，切勿将逻辑联结词混淆。
    \item \textbf{多重集合}：当推广到多个集合时，要注意 $\bigcap$ 和 $\bigcup$ 的变化规律，避免写错。
\end{itemize}
\end{trapbox}